\documentclass[11pt,letterpaper]{article}
\usepackage[margin=0.85in]{geometry}
\usepackage[T1]{fontenc}
\usepackage{graphicx,booktabs,tabularx,longtable,amsmath,xcolor,listings,microtype}
\usepackage[font=small,labelfont=bf]{caption}
\usepackage[section]{placeins}
\usepackage{hyperref}
\hypersetup{colorlinks=true,linkcolor=blue!45!black,citecolor=blue!45!black,urlcolor=blue!45!black,
 pdftitle={ECE 592 Homework 3 Track B: Q1 Q2 Q3 Q4 Q5 Q8 - 105 points},
 pdfauthor={Varunaditya Venkatesan}}
\IfFileExists{report/figures/shared-timing.pdf}{\newcommand{\assetroot}{report/}}{\newcommand{\assetroot}{}}
\newcommand{\repo}{https://github.com/vvenka25/ECE-592-Project-3}
\newcommand{\overleafproject}{https://www.overleaf.com/project/6aeeb80e5d4ccedd231a327b}
\newcommand{\artifact}[1]{\path{#1}}
\newcommand{\rfigure}[4]{\begin{figure}[!htbp]\centering\includegraphics[width=#2\linewidth]{\assetroot figures/#1.pdf}\caption{#3}\label{#4}\end{figure}}
\newcommand{\rtable}[3]{\begin{table}[!htbp]\centering\small\setlength{\tabcolsep}{4pt}\input{\assetroot generated/#1.tex}\caption{#2}\label{#3}\end{table}}
\newcommand{\snippet}[2]{\lstinputlisting[caption={#2}]{\assetroot snippets/#1}}
\newcommand{\evidence}[1]{\par\smallskip\begingroup\small\raggedright\noindent\textit{Evidence and reconstruction.} #1\par\endgroup}
\lstset{basicstyle=\ttfamily\footnotesize,columns=fullflexible,keepspaces=true,
 frame=single,rulecolor=\color{black!25},backgroundcolor=\color{black!2},
 breaklines=true,breakatwhitespace=false,showstringspaces=false,
 aboveskip=6pt,belowskip=6pt,captionpos=b}
\setlength{\parindent}{0pt}
\setlength{\parskip}{4pt plus 1pt}
\widowpenalty=10000\clubpenalty=10000
\renewcommand{\arraystretch}{1.12}
\setlength{\textfloatsep}{10pt plus 2pt minus 2pt}
\setlength{\floatsep}{9pt plus 2pt minus 2pt}
\setlength{\intextsep}{9pt plus 2pt minus 2pt}
\begin{document}
\begin{center}
{\LARGE\bfseries Pipelines, Speculation, and Address Translation\par}
\vspace{5pt}
{\large Microarchitecture Experiments on Two Intel Generations\par}
\vspace{12pt}
{\large Varunaditya Venkatesan\par}
Unity ID: \texttt{vvenka25}\\
ECE 592 --- Homework 3, Track B\\
North Carolina State University, Fall 2026\\
November 7, 2026
\vspace{9pt}

\textbf{GitHub source and evidence:}\\
\url{https://github.com/vvenka25/ECE-592-Project-3}\\[3pt]
\textbf{Overleaf project:}\\
\url{https://www.overleaf.com/project/6aeeb80e5d4ccedd231a327b}
\end{center}
\vspace{2pt}\hrule\vspace{6pt}
\textbf{Selected exercises: Q1, Q2, Q3, Q4, Q5, Q8 --- 105 points.}
\begin{center}\small
\begin{tabular}{clrc lr}
\toprule
Question & Experiment & Points & Question & Experiment & Points\\
\midrule
Q1 & Pipelining & 10 & Q4 & ROB capacity & 30\\
Q2 & Out-of-order execution & 10 & Q5 & Integer physical registers & 35\\
Q3 & Branch prediction & 10 & Q8 & TLB-miss cost & 10\\
\bottomrule
\end{tabular}
\end{center}
\textbf{Abstract.}
We examine six microarchitectural properties on Charnwood's Core i7-6700
(Skylake) and Artemisia's Xeon Gold 5420+ (Sapphire Rapids), with paired
hardware-debug SGX measurements on Artemisia. Fixed-work integer multiplication
shows approximately three active core cycles of dependency latency and a
one-cycle initiation interval on both CPUs. A blocked-older/ready-younger
arithmetic experiment, with dependency and fence controls, supports out-of-order
execution. Balanced shuffled branch outcomes cost substantially more than
regular outcomes, while a branchless control remains nearly unchanged.
Repeated ROB-spacing sweeps place effective overlap loss at N=220--224 on
Charnwood and N=492--500 on Artemisia. Scaled-index LEA writes expose earlier
integer-GPR destination-pressure regions, approximately N=128--144 and
N=200--216; the evidence does not establish total physical register-file size.
For native L1 data-TLB misses served by the STLB, cache/event-controlled
baseline subtraction gives about nine and seven additional active core cycles,
respectively. The enclave retains the principal dependency, branch and
capacity trends, but its protected TLB-event coverage is unqualified. All
results use retained raw measurements; process repetitions, timing units,
controls and attribution limits remain explicit.
\vspace{5pt}\hrule
\clearpage

\section{Experimental setup and common measurements}\label{sec:setup}
\subsection{Machines and placement}
Campaign inventories identify the two CPUs and collection-time configuration
(Table~\ref{tab:machines}). Every selected question has native measurements
on both course machines, as required~\cite{assignment}.
\begin{table}[!htbp]\centering\small
\begin{tabularx}{\linewidth}{lXX}
\toprule
 & Charnwood & Artemisia\\\midrule
CPU & Intel Core i7-6700 & Intel Xeon Gold 5420+\\
Microarchitecture & Skylake & Sapphire Rapids\\
Topology & 1 socket, 4 cores, 8 logical CPUs & 2 sockets, 56 cores, 112 logical CPUs; CPU 68 offline\\
NUMA domains & 1 & 2\\
New-suite placement & CPU 1, node 0; sibling 5 online & CPU 12, node 0; sibling 68 offline\\
SMT control & Private inherited core-scheduling group & No online sibling; before/after activity gate\\
Compiler & GCC 13.3.0, \texttt{-O2} & GCC 11.5.0, \texttt{-O2}\\
SGX access & No current/legacy SGX device path & Hardware-debug SGX with system PSW runtime\\
\bottomrule
\end{tabularx}
\caption{Recorded configuration. Charnwood's controller is on CPU 0; new-suite controllers use another physical core. These are collection-time settings, so reproduction must recheck topology and affinity.}\label{tab:machines}
\end{table}

Affinity precedes allocation. Native memory experiments request node 0 and
4 KiB backing. The final protocol requires full RSS, no transparent huge pages
and read-only per-page NUMA queries before/after timing; it requests no page
migration~\cite{movepages}. Earlier campaigns retain their recorded checks.
These host audits do not verify physical EPC placement. Charnwood's core group
excludes unrelated userspace co-scheduling, with kernel/interrupt contexts
still possible~\cite{coresched}.

\textbf{Preserved Q4 is a distinct configuration.} Artemisia's completed Q4
runs used CPU 1/node 0, with sibling 57 activity recorded as uncontrolled.
Charnwood Q4 also used CPU 1/node 0; an external private core group was applied,
although the legacy manifest says uncontrolled. Its external-control records
identify the actual group; the legacy helper mask also permits CPU 1. Neither
history is rewritten as equivalent to the newer CPU-12 protocol.

\subsection{Samples, work and statistics}
Each treatment keeps 15 timed samples and two fresh-process repetitions; job
and parameter orders are shuffled with recorded seeds. The randomized pointer
cycles cover the full requested footprint and continue between samples.
Full-ring or full-pattern warm-up and a short kernel warm-up precede timing.
All primary observations remain, including recorded ordinary faults and
context switches. Solid circles denote process 1, dashed squares process 2
unless a caption specifies another control style. Shading or error bars show
observed inclusive p10--p90, not confidence intervals. Tables show separate
process medians as \textbf{R1 / R2}; no host, environment or repetition is pooled.

A sample records host TSC ticks. Q1 divides by $32I$ IMULs; Q2, Q3 and Q8 divide
by $I$ blocks, decisions and loads; Q4 and Q5 divide by $2I$ loads. TSC ticks
are not active core cycles. We neither convert ticks using nominal MHz nor
assume equal tick duration across machines. Native PMU core cycles are labeled
separately. Counter brackets include slightly more host-user work than the
TSC interval and require at least 98\% scheduled coverage for qualification.

\clearpage
\subsection{SGX background and same-machine comparison}
Intel SGX provides enclaves: application regions whose code and data receive
hardware-enforced isolation from ordinary host software. SGX applications
separate trusted enclave code from an untrusted host, using explicit enclave
calls to cross that boundary~\cite{sgx}. Here the enclave allocates and
initializes its private arena and executes the assembly kernel; the host
sets affinity, creates the enclave, invokes initialization/warm-up, reads the
timer and counters, and emits raw data. Artemisia used hardware-debug enclaves
(\texttt{SGX\_DEBUG\_FLAG}, \texttt{DisableDebug=0}), not simulation. Debug
mode permits inspection and is a development configuration~\cite{sgx}.
The user previously confirmed Intel SampleEnclave with
\texttt{SGX\_MODE=HW SGX\_DEBUG=1}; that console output was not saved.
Retained custom-enclave semantic and timing checks establish the benchmark's
own hardware execution separately.

Comparing native and SGX on the \emph{same} Artemisia CPU asks whether the same
instruction stream exposes the same behavior, and whether its timing or
transition changes. Within each paired campaign, the kernel bytes, input
pattern or ring geometry, footprint, seed, filler/chain choices, iteration
count, sample count and CPU/node request match. Both environments use host
LFENCE/RDTSCP/LFENCE boundaries around one invocation. The SGX timed interval
includes ECALL entry/exit and marshalling; enclave creation, allocation and
warm-up are outside. Native arenas use anonymous mmap; protected arenas use
enclave allocation, with Q8 explicitly page-aligned in both. Differences in
runtime transitions, protected memory placement, cache/translation state,
frequency or interruptions are possible explanations, not separately measured
causes of a timing difference.

\begin{table}[!htbp]\centering\small
\begin{tabular}{lrr}
\toprule
Enclave setting & Q1/Q2/Q3/Q5/Q8 suite & Preserved Q4\\\midrule
\texttt{HeapMaxSize} & \texttt{0x18000000} (384 MiB) & \texttt{0x30000000} (768 MiB)\\
\texttt{StackMaxSize} & \texttt{0x40000} (256 KiB) & \texttt{0x40000} (256 KiB)\\
\texttt{TCSNum / TCSPolicy} & 1 / 1 & 1 / 1\\
\texttt{ProdID / ISVSVN / DisableDebug} & 0 / 0 / 0 & 0 / 0 / 0\\
\texttt{MiscSelect / MiscMask} & 0 / \texttt{0xFFFFFFFF} & 0 / \texttt{0xFFFFFFFF}\\
\bottomrule
\end{tabular}
\caption{Configuration taken from the benchmark XML and campaign source snapshots. SDK path: \texttt{/opt/intel/sgxsdk}; retained runtime/PSW packages report version 2.23.100.2.}\label{tab:sgx}
\end{table}
\rfigure{shared-timing}{.94}{The common host timer encloses the native invocation or the ECALL. Protected allocation and physical EPC placement differ from native mmap; equal kernel bytes do not make those memory environments identical.}{fig:timer}
\lstinputlisting{\assetroot snippets/timer.c}

\clearpage
\section{Q1: confirm pipelining (10 points)}\label{sec:q1}
\subsection{Method and observable signal}
A pipelined execution resource can start a new operation before the previous
operation produces its dependent result. We therefore measure dependency
latency and initiation interval with fixed instruction work. Each loop executes
32 explicit \texttt{imul \$3,reg,reg} instructions, distributed over one, two
or four independent 64-bit GPR chains. Odd starting values and modulo-$2^{64}$
arithmetic give checked, observable results; the compiler cannot replace the
assembly IMULs with LEAs. With one chain, every multiply waits for its own prior
result. Additional chains make ready multiplies available without changing
the total multiply count.

Each sample contains 100,000 loops (3.2 million IMULs). Register initialization,
checksum generation and invocation are included in the call, so zero-work
controls measure their fixed contribution. We expect per-IMUL time to fall
until execution throughput limits further improvement. That expectation alone
does not distinguish pipelining from replicated execution units.
\snippet{q1.S}{Actual four-chain loop excerpt from \texttt{tb\_q1\_4}. The shown group occurs eight times per loop; the comment marks omitted repetition.}
\rfigure{q1-method}{.94}{Q1 procedure. The small measured plot uses Artemisia native process 1. The arrows denote dependencies, not an instruction-by-instruction hardware trace.}{fig:q1method}

\clearpage
\subsection{Native results, paired enclave result and inference}
\rfigure{q1-results}{.98}{Q1 fixed-work results. Left: native active core cycles per IMUL on both machines. Right: Artemisia host TSC ticks per IMUL in native and hardware SGX. Process repetitions and observed p10--p90 remain separate.}{fig:q1results}
\rtable{q1_native}{Native Q1 process medians in active core cycles/IMUL (R1 / R2).}{tab:q1native}
\rtable{q1_pair}{Paired Artemisia Q1 process medians in TSC ticks/IMUL (R1 / R2).}{tab:q1pair}

Figure~\ref{fig:q1results} and Table~\ref{tab:q1native} show approximately
3.00, 1.50 and 1.00 active core
cycles/IMUL for one, two and four chains, respectively. The one-chain
latency-dominated cost is about three times the sufficiently independent
initiation interval. Zero-work medians are 76 ticks/invocation on Charnwood,
64 natively on Artemisia and 11,046 in its enclave (process 1); no overhead
subtraction is applied. That fixed cost is small per operation in a
3.2-million-IMUL batch.

\textbf{Alternative explanation.} Multiple multiplier units could also speed
independent work. As a separately labeled reference model, LLVM assigns the
selected 64-bit immediate/register IMUL to a single eligible port-1 resource
on Skylake and Sapphire Rapids~\cite{llvm-skl,llvm-spr}. Combined with that
resource bound, the measured latency/throughput gap supports pipelined
multiplication. We have not experimentally counted the units or measured
whole-processor pipeline depth; the model is supporting evidence rather than
a measured specification.

Artemisia SGX preserves the chain-count trend. Its higher ticks coexist with
larger included invocation costs; neither protected arithmetic latency nor a
change in multiplier design is isolated, since enclave PMU coverage is unqualified.
\evidence{\artifact{results/q1-charnwood-private/} and
\artifact{results/q1-artemisia/}: 120 and 240 raw observations, including
zero-work controls. Source: \artifact{trackb/gen_kernels.py}; reconstruction:
\artifact{report/make_report_artifacts.py}.}

\clearpage
\section{Q2: confirm out-of-order execution (10 points)}\label{sec:q2}
\subsection{Method and observable signal}
We place an \emph{older unready arithmetic instruction} before younger ready
work. A pointer load feeds an older IMUL; a later chain of K IMULs uses R9 and
a ready XOR source. Program-order issue would have to wait at the older
unready multiply. Out-of-order issue can execute the younger independent
chain while the load and older multiply are pending.

The dependency control changes the XOR source from ready R11 to the older
result in RAX. A third variant inserts LFENCE before the younger work.
\emph{Every} variant also has a common LFENCE at the block end, so younger
arithmetic cannot hide behind the next iteration's load. This execution
boundary uses Intel's LFENCE ordering semantics~\cite{intel-isa}. The same
random pointer traversal runs with 16 KiB hot and 256 MiB cold footprints.
K=0 exposes baseline/fence cost; K=64 supplies substantial younger work.
Each sample measures 100,000 blocks.
\snippet{q2.S}{Actual overlap kernel \texttt{tb\_q2\_overlap\_64}. The dependent variant reads RAX in the XOR; the fenced variant inserts LFENCE after the older IMUL.}
\rfigure{q2-method}{.94}{Q2 tests bypass of an older blocked operation, beyond independent-stream throughput. The measured response is Artemisia native process 1 at cold K=64.}{fig:q2method}

\clearpage
\subsection{Native results, paired enclave result and inference}
\rfigure{q2-results}{.98}{Q2 cold-data comparisons. Dependence/fence variants are named in the legend; control styles differ from the general process-style convention. Only measured K=0 and K=64 conditions are joined; no intermediate measurements are implied.}{fig:q2results}
\rtable{q2}{Q2 medians in TSC ticks/block, R1 / R2. Hot is 16 KiB; cold is 256 MiB.}{tab:q2}

In Figure~\ref{fig:q2results} and Table~\ref{tab:q2}, cold K=64 independent
younger work takes about 360 ticks/block on
Charnwood, versus approximately 600 with an explicit dependency or preceding
fence. Artemisia native takes about 293--295 ticks independently, versus
430--432 dependent and 439--440 fenced. The matched code order, dependency
change and common block boundary support younger ready arithmetic executing
ahead of the older load-dependent multiply. A speedup from independent
operations alone would not establish this ordering property.

Hot independent K=64 is substantially shorter than cold K=64 on both hosts;
the older cold load provides an appreciable waiting interval. Artemisia's
cold K=0 is close to independent K=64, consistent with hidden arithmetic.
Charnwood's cold K=0 instead exceeds its independent K=64 time. Memory/frequency
state and fence behavior therefore preclude an exact additive latency
identity. The conclusion rests on matched K=64 dependencies and the actual
assembly, not subtraction of an assumed load-plus-compute sum.

Artemisia SGX shows the same ordering of the K=64 controls, with larger cold
times and more variation. This supports the same qualitative behavior; it
does not attribute the extra time to one SGX mechanism. An earlier pilot
without the common block-end fence is retained as a rejected cross-iteration
assay, and the partial first Artemisia collection failed its placement gate.
Neither is pooled with the accepted campaign.
\evidence{\artifact{results/q2-charnwood-private/} (270 rows) and
\artifact{results/q2-artemisia-retry/} (540 rows). Source and reconstruction
are \artifact{trackb/gen_kernels.py} and \artifact{report/make_report_artifacts.py}.}

\clearpage
\section{Q3: confirm branch prediction (10 points)}\label{sec:q3}
\subsection{Method and observable signal}
The conditional kernel reads one byte per decision from the same sequential
64 KiB array. The byte selects one of two equal-length assembly arms; each
adds to an observable outcome count and executes one unconditional jump.
The treatments are constant zero, constant one, alternating, and shuffled
balanced 50/50 outcomes. Seeds 42 and 271828 provide separate shuffled
patterns. Array addresses, footprint and iteration count stay fixed.

A branchless kernel reads the same patterns and computes the same counts with
arithmetic. Its instruction stream is different, so it is a sensitivity
control for outcome-dependent input effects rather than an identical-work
baseline. The branch kernel is explicit assembly, and linked-code checks
verify its conditional jump, arm lengths and absence of if-conversion.
Each timed sample executes 1,048,576 decisions after full-pattern warm-up.
A prediction-sensitive contrast should increase shuffled timing and native
branch misses while leaving the branchless contrast small.
\snippet{q3.S}{Actual conditional hot-loop excerpt from \texttt{tb\_q3\_branch}; both outcome arms remain in the linked native and enclave instruction stream.}
\rfigure{q3-method}{.94}{Q3's pattern and control. The shown shuffled prefix is generated with the recorded seed and 64 KiB shuffle procedure. The small response plot uses Artemisia native process 1; BL denotes the branchless kernel.}{fig:q3method}

\clearpage
\subsection{Native results, paired enclave result and inference}
\rfigure{q3-results}{.98}{Q3 native core-cycle contrasts and Artemisia paired elapsed timings. The native comparison uses core cycles because Charnwood TSC costs drift between processes. All raw timing observations remain retained.}{fig:q3results}
\rtable{q3_native}{Native Q3 medians in active core cycles/decision, R1 / R2. The random seed is part of each treatment label.}{tab:q3native}
\rtable{q3_pair}{Paired Artemisia Q3 medians in TSC ticks/decision, R1 / R2.}{tab:q3pair}

Figure~\ref{fig:q3results} and Table~\ref{tab:q3native} show both native CPUs
predicting the regular stream successfully and incurring much
higher cost for shuffled outcomes. Native shuffled misses are approximately
0.496--0.497 per decision on Charnwood and 0.469--0.473 on Artemisia, compared
with only a few parts per million for alternating outcomes. The branchless
pattern contrast remains small. Combined with the retained conditional jump,
these observations support branch prediction; they do not measure a
calibrated per-misprediction penalty or identify a predictor architecture.

Charnwood alternating conditional TSC medians drift from 5.5566 to 7.3413,
and branchless alternating from 2.5604 to 3.9602, despite stable core-cycle
costs. Elapsed timing alone would confound the comparison. SGX retains the
regular/shuffled timing contrast, but counts only about 0.0003 branches/decision
versus roughly 3 natively: protected miss-rate coverage is insufficient.
\evidence{\artifact{results/q3-charnwood-private/} (240 rows) and
\artifact{results/q3-artemisia/} (480 rows). Generator and raw input hashes
are in \artifact{report/generated/artifact-audit.json}.}

\clearpage
\section{Q4: ROB capacity through overlap loss (30 points)}\label{sec:q4}
\subsection{Preserved method and observable signal}
The completed Q4 implementation, observations, original plots and analyses
are preserved. Each loop loads from two disjoint randomized pointer rings,
with N single-byte NOPs between load A and load B and another N before the
next load A. When the two dependency chains overlap, independent traversal is
cheaper than a serialized traversal of the same complete line set. A long
older miss prevents retirement; increasing filler spacing eventually prevents
useful overlap across the finite instruction window.

The dependent control follows one cycle through all lines. Footprint is
256 MiB, stride 64 bytes, seed 42, one-byte NOPs and loop alignment 0.
Each sample has 100,000 loops and two loads/loop. Charnwood's coarse sweep
selected N=192..240 in steps of four; Artemisia's preserved fine range is
448..576. Each fine run has a fresh order-shuffled repeat. The 2N fillers,
loads and loop control must not be conflated with an established total ROB
entry count.
\snippet{q4.S}{Actual unchanged \texttt{rob\_independent\_224\_b1\_a0} loop, with repeated NOP text explicitly omitted. Both native and enclave link the same kernel object.}
\rfigure{q4-method}{.94}{Q4's two-load spacing assay and serialized control. The small ratio plot is preserved Artemisia native process 1, reconstructed from raw observations without changing their original analysis.}{fig:q4method}

\clearpage
\subsection{Two-machine result and paired native/enclave result}
\rfigure{q4-results}{.98}{Q4 fine observations. Charnwood native shows independent and dependent traversals; Artemisia shows independent native/SGX traversals over a labeled zoom of the preserved range. Orange bands locate the reported operational transitions. The N=548 SGX-repeat peak is retained.}{fig:q4results}
\rtable{q4}{Independent-load medians in TSC ticks/load, R1 / R2. Different N columns reflect each host's own fine bracket, not matched elapsed durations across CPUs.}{tab:q4}

In Figure~\ref{fig:q4results}, Charnwood's largest fine increase is N=220$\rightarrow$224 in both processes,
also under minimum and unflagged-median diagnostics. At N=224,
independent/dependent median ratios are about 0.974 and 1.000. Artemisia's
preserved transition spans N=492--500, with native and hardware-SGX overlap
loss in the same region. These are effective instruction-window bounds for
the actual workload. They demonstrate a larger operational window on
Artemisia, but exact physical ROB-entry attribution remains unproved.

The enclave's higher absolute ticks coexist with a similar principal
transition; the evidence supports a shared qualitative window rather than
a calibrated SGX cost per ROB entry. The isolated N=548 repeat peak would
mislead a global largest-jump-only estimator. Both processes and all primary
samples remain visible, so the transition is not sharpened by trimming that
peak. Uncontrolled SMT on preserved Artemisia Q4 and Charnwood's legacy helper
placement qualify the comparison.

No Artemisia Q4 rerun was needed. The user-confirmed missing Charnwood
comparison was collected from unchanged source/build ID
\texttt{bcd9d621feed36f6}. Its original 3,084-kernel checks and independent
compiled native timer inspection passed. Charnwood retains 1,950 coarse rows
and 390/390 fine rows; the Artemisia fine campaigns retain 1,980 rows each.
\evidence{\artifact{results/q4-coarse-charnwood/},
\artifact{results/q4-fine-charnwood/}, \artifact{results/q4-fine-repeat-charnwood/},
\artifact{results/v2-fine/} and \artifact{results/v2-fine-repeat/}.
The preserved \artifact{ROB_BENCHMARK.md} and
\artifact{ROB_RESULTS_ARTEMISIA.md} remain unchanged; new figure reconstruction
is separate under \artifact{report/}.}

\clearpage
\section{Q5: integer physical-register allocation pressure (35 points)}\label{sec:q5}
\subsection{Register class, method and observable signal}
The tested class is \textbf{64-bit integer general-purpose registers (GPRs)}.
We extend the two-ring layout by replacing NOP filler with
\texttt{lea (reg,reg,2),reg}, which computes $3r$ modulo $2^{64}$ without a
FLAGS destination. The filler avoids plain moves, zero idioms and pure
constant-offset updates. Rotating one, four or eight accumulators changes
dependency depth at the same destination-write count. All variants initialize
the same eight GPR accumulators and preserve the same callee-saved registers.

A four-byte NOP control matches both filler count and encoded length; a
serialized-load LEA control traverses the full same line set. While an older
load prevents retirement, speculative destination writes create pressure on
allocation and execution resources. An early multi-chain transition relative
to the NOP window is evidence for integer-destination pressure; dependence and
scheduler/throughput constraints remain competing limits. Each sample has
10,000 loops. Coarse N=0..512, step 16, selected fine N=120..160 on Charnwood
and 192..232 on Artemisia, step four.
\snippet{q5.S}{Actual eight-chain scaled-index LEA kernel excerpt; repeated filler is marked as omitted. The eight accumulators are R8--R11, RBX, R12, R14 and R15.}
\rfigure{q5-method}{.94}{Q5's destination-pressure treatment. The small measured plot is Artemisia native process 1. The NOP control helps distinguish an earlier destination-sensitive limit from the later instruction-window limit.}{fig:q5method}

\clearpage
\subsection{Native bounds, enclave comparison and conversion limit}
\rfigure{q5-results}{.99}{Q5 fine ranges chosen independently from each host's coarse observations. Native and SGX panels retain four/eight-chain, NOP and serialized-load controls. Process medians and observed p10--p90 remain separate; panel scales differ.}{fig:q5results}
\rtable{q5_native}{Q5 native process medians in TSC ticks/load. N is writes or NOPs per gap; R1 / R2 are separate processes.}{tab:q5native}
\rtable{q5_pair}{Paired Artemisia Q5 process medians in TSC ticks/load, R1 / R2.}{tab:q5pair}

Figure~\ref{fig:q5results} shows Charnwood's four/eight-chain curves rising across roughly N=128--144 while NOPs
remain low; Artemisia's corresponding region is approximately N=200--216.
The NOP transitions occur much later, around coarse N=208--224 and 496--512,
respectively. The direction survives the four/eight-chain change, supporting
an integer-destination-sensitive operational bound. The onset shifts with
chain count, and Charnwood's later LEA time can exceed its serialized control:
allocation, dependency and scheduling/throughput effects are coupled.

Artemisia SGX has the same broad earlier region with greater variation,
including the retained eight-chain process-1 N=216 value of about 403
 ticks/load. Shared bytes do not ensure identical initial physical-register
occupancy in the two environments.

\clearpage
\subsection{Control separation and scope of the register estimate}
\rfigure{q5-coarse-controls}{.98}{The complete Q5 native coarse sweeps retain the one-chain execution diagnostic, multi-chain treatments, serialized loads and NOP control. Green bands mark the multi-chain pressure region; gray bands mark the later NOP-window region. These are actual separate process medians, not schematic capacity curves.}{fig:q5coarse}

Figure~\ref{fig:q5coarse} shows why the fine interval was chosen from Q5's own
coarse data instead of reusing the Q4 NOP bracket. Integer-destination work
changes behavior substantially before NOPs exhaust the instruction window.
The one-chain curve rises early because a long arithmetic dependency itself
costs time; it is not a capacity estimate. Persistence across four and eight
chains is stronger evidence, but execution resources and the scheduler can
still limit how much of the older miss is hidden. The full-range controls
therefore support the destination-sensitive bound without uniquely assigning
it to a physical structure.

The first constant-offset \texttt{lea 1(reg),reg} campaign is retained
separately: its checked results and unexpectedly short one-chain cycles
invalidated the assumed ordinary allocation/execution model. Its mechanism
was not established, and those measurements are not pooled with the revised
scaled-index treatment.

\textbf{What quantity is established.} Conditional on allocation being the
bottleneck, N=128--144 and N=200--216 bracket an effective speculative
integer-destination budget. Precise PRF-versus-scheduler attribution and the
conversion to \emph{total} physical PRF size are unresolved. A conversion
would require measured or justified committed/live mappings, other in-flight
destinations, reclamation timing and reserved capacity. Neither $2N$ nor
$N+16$ supplies that accounting. The experiment is complete; its supported
answer is a conditional bound rather than an invented total size.
\evidence{\artifact{results/q5-charnwood-private/},
\artifact{results/q5-fine-charnwood-private/},
\artifact{results/q5-scaled-residency-artemisia/} and
\artifact{results/q5-fine-artemisia/}. The coarse/fine rows are 4,950/1,650
on Charnwood and 9,900/3,300 on Artemisia.}

\clearpage
\section{Q8: additional cost of an STLB-served TLB miss (10 points)}\label{sec:q8}
\subsection{Miss class, baseline and controls}
The target is an \textbf{L1 data-TLB miss served by the STLB}, rather than a
page walk or an unspecified translation/cache mixture. Both layouts use the
same node count, cache-line count, seed and random dependent traversal order.
Compact nodes occupy consecutive 64-byte lines. Scattered node $i$ occupies
\texttt{base + i*4096 + (i\%64)*64}: one line per 4 KiB page, spread over all
64 L1 set offsets. Both native and enclave arenas are page-aligned; hardware
checks verify actual pointer geometry and cursors.

The 128/256/384 nodes use only 8/16/24 KiB of data in either layout. Full-ring
warming helps keep data hot while scattered translations exceed the L1-TLB
working set. The compact layout is the translation-resident baseline.
Dependency prevents overlap of successive loads. Four simultaneous native
counters record STLB services, completed 4 KiB walks, L1D misses and active
core cycles, using Intel's model-specific event definitions~\cite{perf-skl,perf-spr}.
\snippet{q8.S}{Actual dependent loop \texttt{tb\_q8\_chase}; address construction and warm-up are outside this timed loop.}
\snippet{q8-address.c}{Actual layout arithmetic from \texttt{trackb/engine.c}. Compact and scattered layouts share the same randomized node order.}
\rfigure{q8-method}{.94}{Q8 changes the translation working set while matching touched data lines. Address boxes are schematic; the small cost plot uses the actual class-qualified native process-1 estimates.}{fig:q8method}

\clearpage
\subsection{Qualified native cost and paired enclave observations}
\rfigure{q8-results}{.98}{Left: native cost from separate process medians; no confidence interval is assigned to this derived ratio. Right: Artemisia compact (dotted) and scattered timing. Protected miss attribution remains unqualified.}{fig:q8results}
\rtable{q8_cost}{Additional active core cycles per added STLB service, R1 / R2.}{tab:q8cost}
\rtable{q8_pair}{Artemisia raw layout medians in TSC ticks/load, R1 / R2.}{tab:q8pair}

For Figure~\ref{fig:q8results} and Table~\ref{tab:q8cost}, each matched
process uses the estimate
\[
 C_{\rm extra}=\frac{\operatorname{median}(c_{\rm scattered}/I)-
 \operatorname{median}(c_{\rm compact}/I)}
 {\operatorname{median}(s_{\rm scattered}/I)-
 \operatorname{median}(s_{\rm compact}/I)},
\]
where $c$ counts active core cycles and $s$ counts STLB services. Require
native coverage $\geq98\%$, added services $>0.5$/load, compact services
$<0.05$/load, and both walks and L1D misses $<0.01$/load in either layout.
Both hosts pass: recorded 4 KiB walks are zero, and maximum L1D misses/load
are 0.00273 on Charnwood and 0.0008 on Artemisia. The supported extra costs
are therefore about \textbf{9 and 7 active core cycles}, above the compact
baseline. Page-walk costs are not measured; counter semantics limit precision.

TSC cost also requires effective cycle/TSC medians matching within 2\% and
each condition's p90/p10 ratio $\leq1.02$. All six Artemisia native pairs
pass (about 5.20 ticks/service); four of six Charnwood TSC estimates are
withheld for drift/mismatch. Earlier shifted baselines remain documented.
SGX layouts differ in timing, but unqualified protected counters limit that
result to a translation-working-set contrast rather than a verified miss cost.
\evidence{\artifact{results/q8-clock-charnwood-private/} (180 rows) and
\artifact{results/q8-clock-artemisia/} (360 rows). Final derived quantities and
qualification text are generated in \artifact{report/generated/q8-contrasts.csv}.}

\clearpage
\section{Limitations, reproducibility and disclosure}\label{sec:reproduction}
\subsection{What the evidence does and does not resolve}
All six selected experiments have completed repeated native campaigns on both
machines and corresponding Artemisia hardware-enclave campaigns. Completion
of collection does not make every inferred physical quantity exact. The
principal inference limits are specific:
\begin{itemize}
\item Q1's execution-unit resource bound comes from a cited model, not an experimental unit count; Q2's dependency/fence evidence supports ordering without instruction tracing.
\item Q4 is an operational overlap-loss window, with uncontrolled SMT in its preserved Artemisia runs. Q5 identifies integer GPR pressure, with unresolved precise PRF attribution and total-size conversion.
\item Q8's native class/cache/cycle criteria support its specific STLB-served miss cost. Protected event visibility is insufficient to transfer that attribution to SGX. Full TLB hierarchy reconstruction is not selected.
\item PMU brackets differ slightly from TSC brackets. Frequency, kernel/IRQ/AEX activity, physical EPC placement and other cores' shared-memory traffic remain incompletely controlled. Context-switch/fault flags are observations, not complete interrupt tracing.
\end{itemize}

The failed or rejected paths remain visible: the early cross-iteration Q2
prototype; stopped NUMA-accounting/SMT attempts; the constant-offset LEA
allocation diagnostic; and the three-counter Q8 compact-baseline shift.
They explain methodological choices without replacing the final evidence.
The only requested setup artifact not located is the original SampleEnclave
console log; the user confirmed it was not saved. Its success is explicitly
identified as user-confirmed history, while custom hardware validation logs
and benchmark runtime status are retained. No missing result is filled from
published specifications.

\subsection{Evidence map and reproduction}
The repository is \url{https://github.com/vvenka25/ECE-592-Project-3}.
Its \artifact{REPRODUCING.md} describes the hash-checked restoration of
\artifact{data/raw-results.tar.gz}, which restores recorded manifests, raw
CSVs, mapping audits and source snapshots under \artifact{results/}.
Campaign manifests give commands, parameters, process seeds, CPU/node requests,
source/binary hashes, failures and disturbances. Each question's evidence
paragraph names the actual campaign. The report generator's input hashes and
code-excerpt provenance are in \artifact{report/generated/artifact-audit.json}
and \artifact{report/generated/source-excerpts.json}. No Project 1 content is
an input to the figure/table generator.

From the project root, restore data if using the public checkout and rebuild
this report's numerical tables, vector figures and PDF:
\begin{lstlisting}
python3 data/unpack_results.py   # public checkout only; local data already present
python3 -m pip install -r requirements.txt
MPLCONFIGDIR=/tmp/hw3-mpl python3 report/make_report_artifacts.py
pdflatex -interaction=nonstopmode -halt-on-error -output-directory=report report/report.tex
pdflatex -interaction=nonstopmode -halt-on-error -output-directory=report report/report.tex
\end{lstlisting}
The first command is not needed in the original workspace, where the raw
files are already present. The local dependency is Matplotlib 3.9.4; plain
pdfLaTeX builds the supplied figures and tables without executing Python.
Historical analysis revisions remain alongside existing plots.

\clearpage
Fresh collection is distinct from regenerating figures. Recheck allowed CPUs,
NUMA nodes and SMT state; use fresh output directories. These examples show
new-suite validation and a planned Q3 run, not a request to repeat the completed
campaigns:
\begin{lstlisting}
make -C trackb native
TEST_CPU=12 make -C trackb check
source /opt/intel/sgxsdk/environment
make -C trackb native enclave
TEST_CPU=12 TEST_SGX=1 make -C trackb check
python3 trackb/run.py --question Q3 --cpu 12 --numa-node 0 \
  --environment both --require-idle-sibling --out results/new-q3
# Inspect the plan, then add --execute to collect new observations.
\end{lstlisting}
On Charnwood, the collection helper
\artifact{report/collect_charnwood.sh} records private core scheduling and
controller isolation. Follow Q3$\rightarrow$Q1$\rightarrow$Q2$\rightarrow$Q5$\rightarrow$Q8
for a new collection, and choose Q5 fine bounds from that host's own coarse
observations. Use the isolated
\artifact{report/reproduce_q4_charnwood.py} for unchanged Q4 source; it refuses
execution on Artemisia. Exact historical builds use each campaign's own
source snapshot and toolchain. Current-source reruns are separate evidence.

The generated suite contains 1,053 functions; retained validation checks
semantics, shared native/enclave bytes, native/enclave-host timer assembly,
Q8 page geometry, denied-PMU behavior and native page queries. Native
functional checks do not imply hardware SGX execution: the latter requires
\texttt{TEST\_SGX=1} and is documented in
\artifact{results/trackb-four-counter-hardware-check.log}. Q4's preserved
3,084-function checks and Charnwood timer audit are separate.

\subsection{Overleaf, submission files and AI-use disclosure}
The user-designated project is
\url{https://www.overleaf.com/project/6aeeb80e5d4ccedd231a327b}.
The accompanying Overleaf source ZIP includes \texttt{main.tex}, generated
LaTeX tables, vector figures, source excerpts and build instructions. It
compiles without fetching repository paths or running analysis online.
The normal project URL is included; the private sharing token is omitted.
No successful direct upload is claimed without an authenticated editing
connection. Upload/import of the source ZIP remains a separate publishing
step if such a connection is unavailable.

The assignment requires a ZIP containing the PDF report, source, raw data,
plotting scripts and reproduction instructions, plus an AI-use log for
Track B~\cite{assignment}. The revised submission includes those groups and
the original evidence. The earlier submission and Q4 preservation archive
remain intact; the revised report and new figures are separate artifacts.
GitHub's source-only preservation derivative excludes private signing keys
and generated binaries while retaining the original non-build bytes.

Codex assisted with benchmark/code/test design, hardware collection,
methodological revisions, analysis, figures and report drafting, as recorded
chronologically in \artifact{AI_USE_LOG.md}. For this revision it read the
Project 1 PDF only to guide format, structure and explanation style, read
\artifact{Homework_3.pdf} for requirements, and generated fresh diagrams and
tables from Homework 3 sources and data. It did not copy Project 1 prose,
measurements, code, artwork, conclusions or assignment-specific phases.
No new performance campaign was run for this writing revision. Synthetic
checks are labeled engineering tests, not measurements. The author must
understand and be able to explain the submitted code and conclusions.

\begin{thebibliography}{9}\small\raggedright
\bibitem{assignment} ECE 492/592, \emph{Homework 3: Microarchitecture Reverse Engineering Across Machines}, Fall 2026. Supplied \texttt{Homework\_3.pdf}, pp. 2--4 (Track B selection and common submission requirements).
\bibitem{sgx} Intel, \emph{Intel Software Guard Extensions Web-Based Training}, especially application partitioning and debug enclaves. \url{https://www.intel.com/content/www/us/en/developer/articles/technical/intel-sgx-web-based-training.html}. Accessed November 7, 2026.
\bibitem{llvm-skl} LLVM, \emph{Skylake Client scheduling model}, \texttt{WriteIMul64Imm} resource. \url{https://raw.githubusercontent.com/llvm/llvm-project/main/llvm/lib/Target/X86/X86SchedSkylakeClient.td}. Accessed November 7, 2026.
\bibitem{llvm-spr} LLVM, \emph{Sapphire Rapids scheduling model}, \texttt{WriteIMul64Imm} resource. \url{https://raw.githubusercontent.com/llvm/llvm-project/main/llvm/lib/Target/X86/X86SchedSapphireRapids.td}. Accessed November 7, 2026.
\bibitem{intel-isa} Intel, \emph{Intel 64 and IA-32 Architectures Software Developer's Manual, Instruction Set Reference}, LFENCE. \url{https://www.intel.com/content/www/us/en/developer/articles/technical/intel-sdm.html}.
\bibitem{perf-skl} Intel PerfMon, \emph{Skylake core event definitions}, DTLB load-miss/STLB-hit, completed walks and retired L1D misses. \url{https://raw.githubusercontent.com/intel/perfmon/main/SKL/events/skylake_core.json}. Accessed November 7, 2026.
\bibitem{perf-spr} Intel PerfMon, \emph{Sapphire Rapids core event definitions}. \url{https://raw.githubusercontent.com/intel/perfmon/main/SPR/events/sapphirerapids_core.json}. Accessed November 7, 2026.
\bibitem{coresched} Linux kernel documentation, \emph{Core Scheduling}. \url{https://docs.kernel.org/admin-guide/hw-vuln/core-scheduling.html}.
\bibitem{movepages} Linux man-pages project, \emph{move\_pages(2)}. \url{https://man7.org/linux/man-pages/man2/move_pages.2.html}.
\end{thebibliography}
\end{document}
